% !TeX root = euclid-without-words.tex
% !TeX Program=pdfLaTeX
%
% Diagonals of rhombus bisect each other and are perpendicular
%
\begin{minipage}[t]{.48\textwidth}
\begin{tikzpicture}[scale=.8]
% Bottom side of rhombus
\coordinate (a) at (0,0);
\coordinate (b) at (5,0);
% Draw rhombus at angle 50
\draw[thick,black!10!teal] (a) -- 
  node[left] {$\mathbf{c}$} ++(50:5) coordinate (c) -- 
  node[above] {$\mathbf{c}$} ++(5,0) coordinate (d) -- 
  node[right] {$\mathbf{c}$} (b) -- node[below] {$\mathbf{c}$} cycle;
% Name the two diagonals
\path [name path=d1] (a) -- (d);
\path [name path=d2] (b) -- (c);
% Get their intersection
\path [name intersections={of=d1 and d2,by={intersection}}];
% Draw diagonals
\draw[thick,red] (a) -- node[above] {$\mathbf{a}$} (intersection);
\draw[thick,red] (intersection) -- node[above] {$\mathbf{a}$} (d);
\draw[thick,blue] (b) -- node[right] {$\mathbf{b}$} (intersection);
\draw[thick,blue] (intersection) -- node[right] {$\mathbf{b}$} (c);
% Draw square to denote right angle
\draw[rotate=26] (intersection) rectangle +(8pt,8pt);
\end{tikzpicture}
\end{minipage}
